\chapter{Массивный обмен в слабом взаимодействии, виртуальные частицы и нарушенная симметрия}

Группа Стандартной модели становится гораздо менее загадочной, если понимать её как набор допустимых преобразований в таблице фермионов. Цветовые преобразования идут по вертикали между красным, зелёным и синим. Заряженные слабые преобразования идут по горизонтали внутри кварковых и лептонных дублетов. Электромагнитные фазовые преобразования дают третью операцию. Соответствующие калибровочные бозоны — не украшения, добавленные после того, как найдены симметрии: испускание калибровочного бозона физически осуществляет то же изменение состояния, которое описывает генератор симметрии.

Это наблюдение подводит нас к главному вопросу лекции. В слабой вершине константа связи сравнима по величине с электромагнитной вершиной, и всё же бета-распад колоссально медленнее типичного сильного или электромагнитного процесса. Недостающий ингредиент — распространение. При малом импульсе обмен массивным \(W\) вносит в амплитуду множитель порядка \(1/m_W^2\). Сначала мы разовьём это утверждение в координатном пространстве, затем в пространстве импульсов, проследим его следствия в ходе долгого обсуждения реальных и виртуальных частиц и, наконец, перейдём к первой качественной модели спонтанного нарушения симметрии.

\section{Калибровочная симметрия по таблице фермионов}

\subsection{Цвет действует по вертикали}

Расположим 6 кварковых ароматов в 3 слабых семейства,
\begin{equation}
  (u,d),\qquad(c,s),\qquad(t,b),
\end{equation}
и присвоим каждому кварку красную, зелёную или синюю цветовую метку. Получающаяся кварковая таблица имеет 3 цветовые строки и 6 ароматных столбцов. Цветовая \(SU(3)\) перемешивает 3 элемента в любом одном столбце фиксированного аромата. Например, один из генераторов может превратить красный up-кварк в зелёный up-кварк. На языке частиц та же операция сопровождается испусканием глюона, несущего компенсирующие цвет и антицвет.

Одно и то же цветовое преобразование должно действовать на каждый аромат. Глюон, который превращает красный up-кварк в зелёный up-кварк, — это тот же вид глюона, который превращает красный top-кварк в зелёный top-кварк. Настоящая операция симметрии не может действовать на первое поколение, игнорируя остальные. Следовательно, цветовое преобразование идёт по вертикали и универсально:
\begin{equation}
  q_r\longleftrightarrow q_g\longleftrightarrow q_b,
  \qquad q\in\{u,d,c,s,t,b\}.
\end{equation}
Оно меняет цвет, но не аромат.

Это характерная связь между калибровочной группой и её калибровочными бозонами. Бесконечно малые генераторы описывают преобразования между состояниями, а испускание или поглощение соответствующих бозонов реализует эти преобразования как физические процессы.

\subsection{Слабые преобразования действуют по горизонтали}

Заряженное слабое взаимодействие идёт в другом направлении. Кварк up-типа может испустить заряженный слабый бозон и превратиться в своего down-партнёра, не меняя цвета:
\begin{equation}
  u_b\longrightarrow d_b+W^+.
\end{equation}
В точности тот же \(W^+\) появляется в
\begin{equation}
  c_g\longrightarrow s_g+W^+,
  \qquad
  t_r\longrightarrow b_r+W^+.
\end{equation}
Индексы \(r,g,b\) здесь играют роль спектаторов. Глюоны не переводят состояние по таблице по горизонтали, а заряженные бозоны \(W\) не переводят его по цвету по вертикали.

Добавим электромагнитную фазовую симметрию \(U(1)\), при которой каждое заряженное поле приобретает фазу. Тогда в первом приближении групповая структура такова:
\begin{equation}
  SU(3)\times SU(2)\times U(1).
\end{equation}
Знаки умножения означают, что это независимые множители: один действует на цвет, другой — на слабые дублеты, третий — фазовым поворотом. Это полезная грубая картина, но это ещё не полная история электрослабого взаимодействия.\editorialnote{Точную калибровочную группу Стандартной модели принято записывать как
\(SU(3)_c\times SU(2)_L\times U(1)_Y\). Слабая \(SU(2)_L\) действует только на леворучные дублеты, а окончательная электромагнитная \(U(1)_{\rm em}\) возникает из
\(SU(2)_L\times U(1)_Y\) после нарушения симметрии. В лекции прямо предупреждается, что её первоначальное отождествление двух последних множителей ещё не вполне согласовано; недостающая хиральная и электрослабая структура вводится в последующих лекциях.}

\subsection{Та же схема включает и лептоны}

Под тремя цветовыми строками поместим бесцветную строку лептонов:
\begin{equation}
\begin{array}{c|cc|cc|cc}
 &u&d&c&s&t&b\\ \hline
Q&+\frac23&-\frac13&+\frac23&-\frac13&+\frac23&-\frac13\\
\text{lepton partner}
 &\nu_e&e^-&\nu_\mu&\mu^-&\nu_\tau&\tau^-\\
Q&0&-1&0&-1&0&-1
\end{array}
\end{equation}
В каждой соседней паре электрические заряды отличаются на одну единицу. Именно поэтому один и тот же заряженный слабый бозон может связывать \(u\) с \(d\), \(c\) с \(s\), \(t\) с \(b\), а также заряженный лептон с его нейтрино. Вершине \(W\) важны подходящие квантовые числа и разность зарядов, а не то, является ли фермион кварком или лептоном.

\begin{figure}[htbp]
\centering
\includegraphics[width=0.88\textwidth]{lecture_06_figure_07.png}
\caption{Полная таблица фермионов: сверху три цветовые строки кварков, ниже бесцветная строка лептонов; показаны кварковые заряды и три лептонных дублета.}
\lectureframe{00:18:55}
\end{figure}

Мюон имеет тот же заряд и те же известные внутренние квантовые числа, что и электрон, но примерно в 200 раз тяжелее. Тау-лептон тяжелее ещё больше. Массы нейтрино очень малы и в обсуждении лекции ещё не были известны детально. Пока не найден принцип, который объяснял бы наблюдаемую картину масс семейств.

\begin{classroomqa}
\audiencequestion{Почему лептоны, барионы и мезоны носят такие названия?
\lecturetimestamp{00:19:02}}
\lecturerresponse{Исторически греческие корни означали приблизительно лёгкий, тяжёлый и промежуточный. Эта система названий по массе стала вводить в заблуждение. Мюон когда-то называли мю-мезоном, потому что его масса лежала между массами электрона и протона, но его взаимодействия показывают, что это лептон, гораздо больше похожий на электрон, чем на пион. Сегодня мезон означает кварк--антикварковый адрон и является бозоном; мюон не является мезоном.}
\end{classroomqa}

\section{Загадка слабости}

\subsection{Конкретный процесс с заряженным током}

Возьмём бета-распад нейтрона в качестве рабочего примера. Составы валентных кварков таковы:
\begin{equation}
  n=udd,\qquad p=uud.
\end{equation}
Один down-кварк совершает переход
\begin{equation}
  d\longrightarrow u+W^-,
\end{equation}
а заряженный бозон порождает
\begin{equation}
  W^-\longrightarrow e^-+\bar\nu_e.
\end{equation}
Вместе получаем,
\begin{equation}
  n\longrightarrow p+e^-+\bar\nu_e.
\end{equation}
Два других валентных кварка проходят через процесс как спектаторы. В первой вершине нейтрон превращается в протон; во второй то же слабое взаимодействие, которое связывает лептонный дублет, превращает \(W^-\) в наблюдаемые электрон и электронное антинейтрино.

\subsection{Три константы связи — три очень разных впечатления}

В КЭД заряд электрона \(e\) является амплитудой вершины. Его квадрат даёт характерную вероятность испускания, при этом обычные множители собраны в
\begin{equation}
  \alpha=\frac{e^2}{4\pi}\simeq\frac1{137}
\end{equation}
в натуральных единицах. У сильного взаимодействия есть константа связи \(g_s\), а на типичном адронном масштабе её аналог \(\alpha_s\) может быть порядка \(0.2\). Это гораздо больше, чем \(1/137\), поэтому при ударе по кварку испускание глюонов происходит часто.

Пусть \(g_w\) обозначает константу заряженного слабого взаимодействия. Её квадрат не фантастически мал; на уровне оценок по порядку величины, используемом здесь, он сравним с электромагнитной константой связи. И всё же свободный нейтрон ждёт примерно 15 минут, прежде чем распасться, тогда как сильный распад может происходить за характерное адронное время пролёта. В исторической оценке в лекции используется более старая грубая величина 12 минут, но любая из этих цифр делает одну и ту же загадку неизбежной.\editorialnote{Современное среднее время жизни свободного нейтрона составляет около
\(8.8\times10^2\,{\rm s}\), и между основными методами измерения сохраняется экспериментальное расхождение. Грубое время жизни из лекции сохранено как педагогический масштаб, а не как точное значение.}
\begin{quote}
Если \(g_w\) не чрезвычайно мало, то что делает слабые процессы слабыми?
\end{quote}

Ответ не может заключаться только в вершинах. В каждой диаграмме Фейнмана есть также пропагаторы — множители, описывающие промежуточное распространение от одной вершины к другой.

\section{Массивный пропагатор в координатном пространстве}

\subsection{Вершины — не вся диаграмма}

Диаграмма бета-распада содержит две слабые вершины, вносящие две степени \(g_w\), и линию \(W\), соединяющую их. Обозначим пропагатор в координатном пространстве через \(\Delta(x-y)\), где \(x\) — событие испускания, а \(y\) — событие поглощения или превращения. Это амплитудное ядро, связанное с переносом возмущения поля между этими событиями.

Точный релятивистский пропагатор содержит больше структуры, чем нам нужно. Важная для нас качественная особенность — его зависимость от массы. На малом пространственноподобном расстоянии у него есть степенная сингулярность, схематически
\begin{equation}
  \Delta(d)\sim\frac1{d^2}.
\end{equation}
Безмассовое поле сохраняет степенной хвост. У массивного поля, напротив, возникает экспоненциальное подавление, как только разделение становится большим по сравнению с его комптоновским масштабом. Поэтому график резко загибается вниз.

\subsection{Натуральные единицы задают положение загиба}

Положим
\begin{equation}
  c=\hbar=1.
\end{equation}
Тогда
\begin{equation}
  [m]=[E]=[p]=L^{-1}.
\end{equation}
Знакомые соотношения объясняют каждое равенство: \(E=mc^2\) отождествляет массу и энергию, а \(p\,x\) имеет размерность \(\hbar\), так что импульс и расстояние имеют обратные единицы измерения. Поэтому единственная длина, которую можно построить из массы \(m\), равна
\begin{equation}
  d_m\sim\frac1m.
\end{equation}
Восстанавливая константы, получаем
\begin{equation}
  d_m\sim\frac{\hbar}{mc},
\end{equation}
комптоновскую длину волны. Чем тяжелее обмениваемая частица, тем короче её эффективный радиус действия.

\begin{figure}[htbp]
\centering
\includegraphics[width=0.78\textwidth]{lecture_06_figure_02.png}
\caption{Размерности в натуральных единицах и качественное убывание пропагатора. Чёрная кривая продолжается степенным хвостом; красная кривая для массивного случая загибается вниз вблизи масштаба, обратного массе.}
\lectureframe{00:27:19}
\end{figure}

\begin{figure}[htbp]
\centering
\begin{tikzpicture}[x=1.0cm,y=0.9cm]
  \draw[->] (0,0) -- (7,0) node[right] {$d$};
  \draw[->] (0,0) -- (0,5) node[above] {$|\Delta(d)|$};
  \draw[thick]
    plot[smooth] coordinates {(0.45,4.5) (0.9,3.3) (1.6,2.5)
      (2.5,1.9) (3.7,1.45) (5.1,1.12) (6.4,0.9)};
  \draw[thick,red]
    plot[smooth] coordinates {(0.45,4.5) (0.9,3.3) (1.6,2.5)
      (2.5,1.9) (3.2,1.55) (3.7,0.95) (4.2,0.36) (4.8,0.1)};
  \draw[dashed] (3.4,0) -- (3.4,1.4);
  \node at (5.7,1.4) {$m=0$};
  \node[red] at (4.75,0.72) {$m>0$};
  \node[below] at (3.4,0) {$d\sim m^{-1}$};
\end{tikzpicture}
\caption{Аккуратная схема рассуждения, приведённого на доске. Это не точная формула для четырёхмерного пропагатора; она выделяет диапазон, определяемый массой.}
\end{figure}

\begin{classroomqa}
\audiencequestion{Что отложено по вертикальной оси?
\lecturetimestamp{00:28:30}}
\lecturerresponse{Величина пропагатора, обозначаемая \(\Delta\). Её квадрат вносит вклад в вероятность после учёта полной амплитуды, внешних состояний и факторов фазового пространства. График призван показать, что амплитуда обмена массивной частицей на большом расстоянии быстро подавляется.\editorialnote{Сам по себе пропагатор не является непосредственно измеримой амплитудой вероятности того, что частица движется по классической траектории. Формулировка лекции полезна как оценка радиуса действия: пропагатор представляет собой один из множителей в полной амплитуде рассеяния или распада.}}
\end{classroomqa}

\section{Преобразование Фурье и предел Ферми}

\subsection{Эксперименты определяют импульсы}

В эксперименте обычно измеряют внешние импульсы, а не две пространственно-временные точки, в которых внутренняя линия начинается и заканчивается. При бета-распаде импульс, переносимый \(W\), определяется по разности импульсов нейтрона и протона и равен суммарному импульсу электрона и антинейтрино. Пропагаторы в координатном и импульсном представлениях связаны преобразованием Фурье.
\begin{equation}
  k=p_n-p_p=p_e+p_{\bar\nu_e}.
\end{equation}

В соглашении о положительном знаменателе, использованном на доске,
\begin{equation}
  \widetilde{\Delta}_W(k)=\frac1{k^2+m_W^2},
  \qquad
  \widetilde{\Delta}_\gamma(k)=\frac1{k^2}.
\end{equation}
Пропагатор безмассового фотона становится большим при малой передаче импульса. Пропагатор массивного \(W\) никогда не превышает величины порядка \(1/m_W^2\). Для бета-распада
\begin{equation}
  |k^2|\ll m_W^2,
\end{equation}
поэтому деталями \(k^2\) можно пренебречь, и
\begin{equation}
  \widetilde{\Delta}_W(k)\simeq\frac1{m_W^2}.
\end{equation}

\begin{classroomqa}
\audiencequestion{Должно ли релятивистское выражение для \(k^2\) содержать знак минус между энергией и пространственным импульсом?
\lecturetimestamp{00:35:11}}
\lecturerresponse{Да. При метрике Минковского с преимущественно отрицательной сигнатурой \(k^2=(k^0)^2-\boldsymbol{k}^2\), а скалярный знаменатель по соглашению имеет вид \(k^2-m_W^2+i\epsilon\). На доске используется эквивалентная положительная пространственно-подобная, или евклидова, величина. В низкоэнергетическом бета-распаде передаваемый четырёхимпульс очень мал по сравнению с \(m_W\), так что в обоих обозначениях получается одно и то же подавление \(1/m_W^2\) с точностью до несущественного знака в амплитуде.}
\end{classroomqa}

\subsection{Низкоэнергетическая амплитуда}

Две слабые вершины и один пропагатор дают
\begin{align}
  \mathcal A_{\rm weak}
  &\sim g_w\,\widetilde{\Delta}_W(k)\,g_w\\
  &\sim\frac{g_w^2}{k^2+m_W^2}
  \simeq\frac{g_w^2}{m_W^2}.
\end{align}
Возведение амплитуды в квадрат даёт характерное подавление
\begin{equation}
  |\mathcal A_{\rm weak}|^2
  \sim\frac{g_w^4}{m_W^4},
\end{equation}
до учёта кинематических факторов и факторов фазового пространства, соответствующих данному распаду. Слабое взаимодействие является слабым при ядерных энергиях не потому, что его элементарная калибровочная константа связи нелепо мала, а потому, что его переносчик тяжёлый.

\begin{figure}[htbp]
\centering
\includegraphics[width=0.86\textwidth]{lecture_06_figure_03.png}
\caption{Слабый переход нейтрон--протон и перенос импульса внутренней линией, как обсуждалось при выборе знака пропагатора.}
\lectureframe{00:35:40}
\end{figure}

\begin{figure}[htbp]
\centering
\includegraphics[width=0.84\textwidth]{lecture_06_figure_04.png}
\caption{Топология бета-распада рядом с заключённой в рамку низкоэнергетической оценкой \(g_w^2/m_W^2\).}
\lectureframe{00:37:25}
\end{figure}

Численно \(m_W\) с точностью этой лекции составляет примерно сто масс протона. Множитель \(m_W^{-2}\) в амплитуде превращается в \(m_W^{-4}\) в скорости распада, так что даже константа связи обычного для калибровочной теории размера приводит к очень большому времени жизни.

\begin{classroomqa}
\audiencequestion{Речь идёт о подавлении на уровне одной миллионной или одной миллиардной, и как изначально узнали массу \(W\)?
\lecturetimestamp{00:38:11}}
\lecturerresponse{Грубая оценка в лекции указывает на малость низкоэнергетической слабой вероятности на уровне одной миллиардной по сравнению с неподавленным адронным процессом. Исторически низкоэнергетические эксперименты не измеряли по отдельности \(g_w\) и \(m_W\). Они измеряли малую комбинацию \(g_w^2/m_W^2\), называемую константой Ферми. Когда калибровочную константу связи связали с электрослабой теорией, из этой комбинации следовало существование очень тяжёлого переносчика; позже ускорители непосредственно породили \(W\) и измерили его массу.}
\end{classroomqa}

При импульсах, гораздо меньших, чем \(m_W\), во всех экспериментах проявляется почти одна и та же константа:
\begin{equation}
  G_F\sim\frac{g_w^2}{m_W^2}.
\end{equation}
Именно поэтому Ферми смог описать бета-распад в 1930-х годах как точечное четырёхфермионное взаимодействие задолго до наблюдения переносчика. По мере роста энергий столкновений эксперименты стали чувствительны к импульсной зависимости, скрытой в пропагаторе, что подтвердило: эффективная константа была низкоэнергетической тенью массивного обмена.\editorialnote{В стандартной электрослабой нормировке на уровне дерева,
\[
  \frac{G_F}{\sqrt2}=\frac{g^2}{8m_W^2}.
\]
Лекция намеренно опускает численные и спинорные множители и сохраняет только масштабную зависимость \(G_F\sim g_w^2/m_W^2\). \(W\) и \(Z\) были непосредственно открыты в CERN в 1983 году; упомянутая в лекции датировка серединой и концом 1970-х слишком ранняя для прямого наблюдения.}

Исторически идея о том, что слабые взаимодействия могут опосредоваться тяжёлой частицей, возникла в середине XX века; в лекции среди теоретиков, рассматривавших такую возможность, упоминается Julian Schwinger, и отмечается, что буква \(W\), вероятно, относилась к слабому взаимодействию. Низкоэнергетические данные уже показывали, что такой переносчик должен быть чрезвычайно тяжёлым в ядерном масштабе, задолго до того, как ускорители получили достаточно энергии, чтобы рождать его как реальную частицу.

\subsection{Три независимых фактора, определяющих скорость процесса}

Вопрос из аудитории о ядерном связывании позволяет разделить три логически различные фактора, управляющие переходом:
\begin{enumerate}
  \item константы связи в вершинах;
  \item пропагаторы, включая массы и зависимость от импульса;
  \item доступная энергия при переходе из начального состояния в конечное и фазовое пространство.
\end{enumerate}
Связывание может закрыть канал распада, потому что конечное состояние имело бы большую энергию, чем начальное. Это кинематический запрет, а не дополнительное уменьшение \(g_w\) и не изменение \(m_W\). Если канал остаётся открытым, но доступной энергии мало, его фазовое пространство всё же может быть сильно подавлено.

\begin{classroomqa}
\audiencequestion{Подавляют ли поле или связывание с соседним протоном бета-распад через тот же механизм пропагатора?
\lecturetimestamp{00:41:13}}
\lecturerresponse{Не в обсуждаемом примере. Связывание важно потому, что оно изменяет энергии начального и конечного ядерных состояний. Процесс не может происходить, если конечное состояние энергетически недоступно. Константы связи, пропагаторы и доступная энергия — три отдельных составляющих.}
\end{classroomqa}

\begin{classroomqa}
\audiencequestion{Действительно ли константы связи постоянны?
\lecturetimestamp{00:42:55}}
\lecturerresponse{Только на фиксированном масштабе. Квантовые поправки заставляют эффективные константы связи зависеть от импульса или расстояния. В настоящем рассуждении каждая константа связи берётся на масштабе, релевантном данному процессу; зависимость от масштаба будет рассмотрена позже.}
\end{classroomqa}

\section{Виртуальный обмен без нарушения законов сохранения}

\subsection{Почему может появиться тяжёлая линия}

Обычный распад нейтрона не содержит достаточно энергии, чтобы породить свободный \(W\), лежащий на массовой оболочке. \(W\) на диаграмме является внутренним и находится вне массовой оболочки. В лекции этот факт вводится с помощью традиционной оценки энергия--время
\begin{equation}
  \Delta E\,\Delta t\lesssim\hbar:
\end{equation}
промежуточная конфигурация, для которой классическая энергия как будто отличается на \(\Delta E\), может вносить вклад лишь в течение времени порядка \(\hbar/\Delta E\).

Аналогия с туннелированием делает задуманную картину наглядной. Поместим альфа-частицу в ядерную потенциальную яму. Её энергия лежит ниже вершины барьера, так что классически она не может выйти наружу. С квантово-механической точки зрения волновая функция проникает сквозь барьер, и частица может появиться по ту сторону в точке с той же полной энергией, что и первоначально. Если попытаться описать внутренность барьера как обычную классическую траекторию, потенциальная энергия оказывается слишком большой, а кинетическая энергия должна была бы быть отрицательной. Эта несостоятельная классическая картина аналогична попытке рассматривать внутренний \(W\) как реальную частицу.

\begin{figure}[htbp]
\centering
\includegraphics[width=0.78\textwidth]{lecture_06_figure_06.png}
\caption{Аналогия с туннелированием: начальная точка и отмеченная точка на дальнем склоне имеют одинаковую полную энергию, тогда как промежуточная область барьера классически запрещена.}
\lectureframe{00:46:00}
\end{figure}

\begin{figure}[htbp]
\centering
\begin{tikzpicture}[x=1.0cm,y=0.9cm]
  \draw[->] (0,0) -- (7,0) node[right] {$x$};
  \draw[->] (0,0) -- (0,4.8) node[above] {$V(x)$};
  \draw[thick]
    plot[smooth] coordinates {(0.3,1.4) (1.0,0.7) (1.7,0.55)
      (2.1,0.9) (3.2,3.9) (4.2,2.1) (5.0,0.75) (6.6,0.18)};
  \draw[dashed] (1.7,0.75) -- (5.0,0.75);
  \fill (1.7,0.75) circle (0.06);
  \fill (5.0,0.75) circle (0.06);
  \node at (3.35,0.35) {$\Delta E\,\Delta t\lesssim\hbar$};
\end{tikzpicture}
\caption{Наглядная реконструкция барьера и точек с одинаковой энергией.}
\end{figure}

В картине бета-распада нейтрон может флуктуировать в протон и внутренний \(W^-\), но эта тяжёлая линия должна быстро исчезнуть. Он может снова поглотиться, восстанавливая нейтрон, или закончиться во второй вершине и породить гораздо более лёгкие электрон и антинейтрино. Полные начальная и конечная энергии совпадают.\editorialnote{Энергия и импульс сохраняются в каждой фейнмановской вершине. Виртуальные частицы не заимствуют энергию буквально и не возвращают её позже. Точное утверждение состоит в том, что по внутренней линии интегрируют по четырёхимпульсам, которые не обязаны удовлетворять соотношению массовой оболочки для реальной частицы \(k^2=m_W^2\). Язык принципа неопределённости и туннелирования сохранён, потому что именно таков объяснительный ход лекции, но утверждение о нахождении вне массовой оболочки и есть современная формулировка квантовой теории поля.}

\subsection{Реальные W-бозоны}

Достаточно энергичный ускоритель может породить реальный \(W\). Тогда у него есть собственное время жизни, не связанное с предполагаемым временем неопределенности виртуальной линии. Разогнанный реальный \(W\) живет в лабораторной системе дольше вследствие обычного релятивистского замедления времени. При двухчастичном лептонном распаде покоящегося \(W\)
\begin{equation}
  W^-\longrightarrow e^-+\bar\nu_e,
\end{equation}
электрон и антинейтрино уносят большие кинетические энергии и разлетаются почти в противоположных направлениях; каждая из этих энергий имеет порядок \(m_W/2\). Энергия покоя не исчезла; она превратилась в энергию конечного состояния. Возможны также мюонный и тау-каналы, а также кварковые каналы.

\begin{classroomqa}
\audiencequestion{Сколько живет реальный \(W\), и можно ли вычислить его время жизни по той же константе связи?\lecturetimestamp{00:49:49}}
\lecturerresponse{Да. Вычисляют амплитуду каждого разрешенного распада, возводят ее в квадрат, интегрируют по фазовому объему конечного состояния и суммируют парциальные ширины. Обратная величина полной ширины дает время жизни. Озвученная в расшифровке численная оценка искажена, но качественный смысл состоит в том, что собственное время жизни чрезвычайно мало, а для движущегося \(W\) имеет место обычное замедление времени.\editorialnote{Используя измеренную ширину
\(\Gamma_W\simeq2.1\,{\rm GeV}\), получаем
\(\tau_W=\hbar/\Gamma_W\simeq3.2\times10^{-25}\,{\rm s}\). Это современное значение устраняет искаженный численный фрагмент диалога, не приписывая его лекции.}}
\end{classroomqa}

\begin{classroomqa}
\audiencequestion{Как квадрат амплитуды превращается в скорость распада и, следовательно, во время жизни?\lecturetimestamp{00:50:52}}
\lecturerresponse{Квадрат матричного элемента объединяют с плотностью доступных конечных состояний, дельта-функциями законов сохранения и нормировочными множителями, чтобы получить вероятность в единицу времени, то есть ширину распада \(\Gamma\). Вероятность нераспада равна \(e^{-\Gamma t}\) в единицах, где \(\hbar=1\), так что \(\tau=1/\Gamma\), или \(\tau=\hbar/\Gamma\) после восстановления размерных единиц. Поэтому малая ширина по слабому взаимодействию означает большое время жизни.}
\end{classroomqa}

\begin{classroomqa}
\audiencequestion{Можно ли приписать скорость частице, туннелирующей через барьер?\lecturetimestamp{00:52:42}}
\lecturerresponse{Можно разделить оцененное расстояние на оцененное время прохождения, но эту величину не следует интерпретировать как скорость классической частицы или сигнала. Попытка измерить траекторию внутри барьера изменяет состояние и разрушает исходный эксперимент по туннелированию.}
\end{classroomqa}

\section{Что значит наблюдать виртуальную частицу?}

\subsection{Внешние частицы и внутренние линии}

Реальные частицы — это входящие и выходящие асимптотические состояния, которые можно приготовить или зарегистрировать далеко от области взаимодействия. Виртуальные частицы — это внутренние линии в пертурбативном представлении амплитуды. Это не дополнительные скрытые следы в детекторе.

\begin{classroomqa}
\audiencequestion{Существует ли четкое определение различия между реальной и виртуальной частицей?\lecturetimestamp{00:55:05}}
\lecturerresponse{Реальная частица — это внешнее состояние на массовой оболочке, входящее в область взаимодействия или покидающее ее. Виртуальная частица — это внутренняя линия вне массовой оболочки на диаграмме Фейнмана. Последняя является частью математического разложения, а не независимо наблюдаемым объектом.}
\end{classroomqa}

Чтобы разрешить процесс, длящийся всего \(\Delta t\), зонду нужны частоты порядка \(1/\Delta t\) и, следовательно, энергия порядка \(\hbar/\Delta t\). Фотон, который может с достаточным разрешением различить альфа-частицу под высоким барьером, может сообщить ей энергию, необходимую для подъема в область барьера. Точно так же зонд, достаточно энергичный, чтобы разрешить короткоживущий внутренний \(W\), может создать реальный \(W\). Попытка наблюдения превращает исходный процесс в другой эксперимент.

В лекции это формулируется как постоянно возникающее квантово-механическое затруднение: вмешательство, необходимое для проверки запрещенной промежуточной картины, как раз и сообщает ту энергию, которая делает кажущуюся запрещенной конфигурацию реальной.

\subsection{Флуктуации протона}

В теории возмущений состояние протона содержит амплитуды конфигураций, таких как
\begin{equation}
  p\longleftrightarrow n+W^+,
\end{equation}
а у нейтрона аналогично есть
\begin{equation}
  n\longleftrightarrow p+W^-.
\end{equation}
Если не вмешивается никакое другое взаимодействие, внутренний бозон вновь поглощается. Жесткий зонд может вместо этого оставить \(W\) среди выходящих частиц. Это можно назвать выбиванием флуктуации или рождением \(W\). Формулировки различаются; измеримое утверждение — одна и та же амплитуда конечного состояния.

\begin{classroomqa}
\audiencequestion{Когда жесткий зонд порождает выходящий \(W\), выбивает ли он частицу, которая уже там была, или создает новую?\lecturetimestamp{01:01:20}}
\lecturerresponse{Для организации вычислений можно использовать любую из этих картин. Одно описание с временным упорядочением говорит, что протон флуктуировал в \(n+W^+\), а зонд выбил нейтрон до повторного поглощения; другое говорит, что конечный \(W\) был создан в столкновении. Инвариантной является вычисленная вероятность наблюдаемого конечного состояния, а не некая единственная классическая история о том, что существовало до измерения.}
\end{classroomqa}

Увеличение энергии столкновения не выявляет фиксированный перечень объектов, которые были скрыты внутри мишени. Оно открывает множество возможных конечных состояний: один \(W\), несколько бозонов, кварковые струи, адроны и другие частицы, разрешенные законами сохранения квантовых чисел. Достаточно жесткое столкновение даже с элементарным полем может породить протоны и многие другие составные объекты. То, что выходит, не обязано было заранее присутствовать в виде составляющей. В лекции эта мысль доводится до комического предела: если это допускают энергия и квантовые числа, то даже столы, стулья или атом водорода логически не исключены как фантастически маловероятные конечные конфигурации. Эта шутка содержит серьезное предостережение против того, чтобы определять элементарность по наивному перечню продуктов столкновения.

\subsection{Вакуумные флуктуации и температура}

Флуктуация, связанная с протоном, — не то же самое, что вакуумная диаграмма. Первая является компонентой состояния взаимодействующего протона; вакуумная флуктуация не имеет внешней протонной линии. В обоих случаях достаточно энергичный зонд может породить реальные частицы, соответствующие разрешенному конечному состоянию. Например, два энергичных фотона могут породить электрон--позитронную пару, когда энергия в системе центра масс превышает порог.

\begin{classroomqa}
\audiencequestion{Зависят ли эти реакции от температуры?\lecturetimestamp{01:05:17}}
\lecturerresponse{Элементарной амплитуде не требуется тепловая среда. Температура важна статистически: высокотемпературная система содержит много частиц с большой кинетической энергией, поэтому столкновения с превышением порога становятся обычными. Нагрев системы, содержащей связанную альфа-частицу, тоже может сообщить достаточно энергии для надбарьерного выхода в дополнение к квантовому туннелированию. Несколько изолированных энергичных частиц образуют пучок; множество энергичных частиц с тепловым распределением образуют горячую среду.}
\end{classroomqa}

\begin{classroomqa}
\audiencequestion{Является ли \(W\), испускаемый при флуктуации протона, флуктуацией вакуума?\lecturetimestamp{01:06:56}}
\lecturerresponse{Нет. Внутренний \(W\) в
\(p\leftrightarrow n+W^+\) является флуктуацией состояния взаимодействующего протона. Замкнутый виртуальный процесс без внешнего протона — это вакуумная флуктуация. Они подчиняются одним и тем же правилам теории поля, но относятся к разным амплитудам.}
\end{classroomqa}

\begin{classroomqa}
\audiencequestion{Почему антинейтрино необходимо при лептонной вершине \(W\)?\lecturetimestamp{01:08:51}}
\lecturerresponse{Оно обеспечивает сохраняющиеся лептонные квантовые числа и входит в базовую вершину заряженного тока, связывающую \(e\), \(\nu_e\) и \(W\). Черта указывает, что это вылетающее электронное антинейтрино, а не нейтрино.}
\end{classroomqa}

\section{Насколько далеко можно зайти в превращении энергии?}

Далее на занятии правило о том, что энергия может рождать частицы, доводится до намеренно абсурдного предела. Может ли кофейное зерно превратиться в шоколадный торт? Практического процесса такого рода не существует, но релятивистская квантовая теория не налагает правила, согласно которому конечное вещество должно напоминать исходный объект. Она требует достаточной энергии, правильных сохраняющихся зарядов и ненулевой амплитуды. Килограмм имеет энергию порядка
\begin{equation}
  Mc^2\sim9\times10^{16}\ {\rm J},
\end{equation}
так что любое микроскопическое столкновение, которое собрало бы конкретный объект килограммового масштаба, было бы запредельно маловероятным.

Планковская масса,
\begin{equation}
  m_{\rm P}=\sqrt{\frac{\hbar c}{G}}
  \simeq2.2\times10^{-5}\ {\rm g},
\end{equation}
не является максимально возможной массой. Это масштаб, на котором гравитация становится сильной для кванта, локализованного на своей комптоновской длине волны. В одном килограмме содержится примерно \(5\times10^7\) планковских масс.

\begin{classroomqa}
\audiencequestion{Если столкновения превышают планковский масштаб, могут ли возникать сколь угодно тяжелые объекты, такие как стулья или планеты?\lecturetimestamp{01:10:37}}
\lecturerresponse{При транс-планковской энергии в системе центра масс, сосредоточенной в малой области, подавляюще естественным исходом является гравитационный коллапс: черная дыра. Высокоорганизованный макроскопический объект был бы чрезвычайно редкой флуктуацией. В лекции привлекается обращение времени: поскольку стул может упасть в черную дыру, обратная амплитуда не равна в точности нулю. Ее вероятность подавлена огромным энтропийным фактором, так что возможность не следует смешивать с практической вероятностью.}
\end{classroomqa}

Это отступление подкрепляет сделанное ранее предупреждение. Квантовая теория поля предсказывает амплитуды между заданными начальными и конечными состояниями. Наглядное утверждение, будто каждый продукт буквально содержался в мишени до столкновения, не является частью этого предсказания.

\section{Бозоны, их взаимодействия и их статистика}

\subsection{Фотоны могут рассеиваться}

То, что частица является бозоном, не означает, что она не взаимодействует. Бозон или фермион задаёт статистику обмена; вершины взаимодействия — отдельные динамические факты. Классические световые пучки обычно проходят друг сквозь друга, потому что в КЭД нет четырёхфотонной вершины древесного уровня, а низкоэнергетическое фотон--фотонное рассеяние чрезвычайно мало. Тем не менее два фотона могут взаимодействовать через виртуальную электрон--позитронную петлю:
\begin{equation}
  \gamma+\gamma
  \longrightarrow(e^-e^+)_{\rm virtual}
  \longrightarrow\gamma+\gamma.
\end{equation}
Если энергия в системе центра масс достаточно велика, они могут вместо этого породить реальную пару,
\begin{equation}
  \gamma+\gamma\longrightarrow e^-+e^+.
\end{equation}
Энергии оптических фотонов составляют несколько электрон-вольт, тогда как порог рождения пары имеет порядок \(2m_e c^2\approx1.02\,{\rm MeV}\). Поэтому виртуальный процесс при обычных оптических энергиях сильно подавлен, и именно поэтому пересекающиеся световые пучки, по-видимому, не замечают друг друга.

\begin{classroomqa}
\audiencequestion{Взаимодействуют ли виртуальные \(W\)-бозоны с другими бозонами просто потому, что они бозоны?
\lecturetimestamp{01:12:47}}
\lecturerresponse{Они могут взаимодействовать, но не просто потому, что подчиняются статистике Бозе. Динамические взаимодействия определяются калибровочными вершинами. Статистика же определяет, как ведут себя многочастичные амплитуды при перестановке тождественных частиц. Фотоны слабо взаимодействуют через заряженные виртуальные частицы, а электрослабая теория содержит также прямые самовзаимодействия неабелевых \(W\)-бозонов.}
\end{classroomqa}

\subsection{Бозе-усиление — не сила}

Различие между фермионной и бозонной статистикой становится заметным, когда несколько тождественных частиц занимают близкие квантовые состояния. Бозоны благоприятствуют многократному заселению, а фермионы исключают одинаковое заселение. В типичном столкновении при высоких энергиях частицы разлетаются в хорошо разделённые состояния по импульсу, поэтому этот статистический эффект может быть менее очевиден, чем в лазере или конденсате.

В принципе, \(W\)-бозоны обладают тем же бозе-усилением. В лекции в шутку говорится, что можно было бы вообразить систему, в которой осуществляется лазерная генерация \(W\)-бозонов. Их огромная масса и ультракороткое время жизни делают такое устройство фантастически непрактичным, но статистика у них действительно бозонная.

\begin{classroomqa}
\audiencequestion{Какого цвета протон или нейтрон?
\lecturetimestamp{01:18:41}}
\lecturerresponse{Физический адрон является цветовым синглетом; его часто называют бесцветным или белым. Краткое обозначение говорит, что барион содержит один красный, один зелёный и один синий кварк. Точнее, его цветовая волновая функция представляет собой антисимметричную суперпозицию
\[
  |{\rm singlet}\rangle
  =\frac1{\sqrt6}\epsilon_{abc}|q^a q^b q^c\rangle,
\]
а не классическое приписывание постоянного цвета каждому названному кварку.}
\end{classroomqa}

\begin{classroomqa}
\audiencequestion{Будет ли в этом курсе обсуждаться конденсация Бозе?
\lecturetimestamp{01:19:25}}
\lecturerresponse{Да, через явление Хиггса. Ненулевое поле Хиггса во всём вакууме в важных отношениях аналогично бозонному конденсату. Эта аналогия будет развита после того, как будет введено спонтанное нарушение симметрии.\editorialnote{Вакуумное среднее поля Хиггса не является в буквальном смысле обычным газом хиггсовских частиц на массовой оболочке, занимающих одно состояние. Конденсация — это проясняющая аналогия для когерентного бозонного поля и выбранного вакуума, а не полное микроскопическое отождествление.}}
\end{classroomqa}

\section{Недостающий нейтральный бозон слабого взаимодействия}

Вернёмся теперь к таблице групп. Если горизонтальная слабая симметрия есть \(SU(2)\), то её размерность равна
\begin{equation}
  \dim SU(2)=3
\end{equation}
и потому имеется три калибровочных поля. Заряженные комбинации \(W^+\) и \(W^-\) объясняют лишь два из них. Требуется третье, нейтральное поле. Оно тесно связано с наблюдаемым бозоном \(Z\), но не является просто фотоном. Точные нейтральные комбинации включают и третье поле \(SU(2)_L\), и поле \(U(1)_Y\). Это смешивание вместе с правильным определением электрического заряда намеренно откладывается до следующей лекции.

\section{Спонтанное нарушение симметрии}

\subsection{Симметричный выбор может иметь неустойчивую неопределённость}

Представьте круглый обеденный стол, где у каждой сервировки есть прибор слева и прибор справа. Пока никто не двинулся, левое и правое эквивалентны. Как только один гость тянется вправо, каждый соседний выбор оказывается вынужденным, и весь стол устанавливается в правостороннюю конфигурацию. Правила были симметричны, но реализованная конфигурация выбрала одну из двух возможностей. Это лишь аналогия, однако она выделяет нужное нам сочетание: симметрию, несколько эквивалентных исходов и неустойчивость, позволяющую сколь угодно малому возмущению глобально выбрать один исход.

\subsection{Решётка орёл--решка}

Для более чистой модели поместим в каждый узел решётки переменную \(s_i=\pm1\) с двумя состояниями. Назовём \(+1\) орлом, а \(-1\) решкой. Пусть соседние узлы предпочитают выстраиваться одинаково:
\begin{equation}
  H=-J\sum_{\langle ij\rangle}s_i s_j,
  \qquad J>0.
\end{equation}
Параллельные соседи дают меньшую энергию, чем антипараллельные. Поэтому существуют два вырожденных основных состояния,
\begin{equation}
\begin{aligned}
  s_i&=+1\quad\text{for every }i,
  \\
  s_i&=-1\quad\text{for every }i.
\end{aligned}
\end{equation}
Гамильтониан инвариантен относительно глобального переворота \(s_i\to-s_i\), но каждое из выбранных основных состояний — нет.

Теперь зафиксируем один очень далёкий граничный спин в положении «орёл». В обычной короткодействующей системе мы могли бы ожидать, что одна удалённая степень свободы здесь никак не повлияет. Однако в упорядоченной фазе это граничное условие выбирает вакуум «все орлы». После такого выбора одиночная «решка» в море «орлов» требует энергии, поскольку не согласуется со своими соседями. Локальные возбуждения определяются относительно выбранного основного состояния, хотя лежащие в основе уравнения сохраняют симметрию орёл--решка.

\subsection{Явное нарушение и доменные стенки}

Явное нарушение добавило бы локальное смещение,
\begin{equation}
  H_{\rm explicit}
  =-J\sum_{\langle ij\rangle}s_i s_j-h\sum_i s_i,
  \qquad h\ne0.
\end{equation}
Тогда орёл и решка не имеют одинаковой энергии даже в одном изолированном узле, и существует только один предпочтительный вакуум. При спонтанном нарушении \(h=0\): локальные правила симметричны и два вакуума остаются вырожденными, но система занимает один из них.

Потребуем вместо этого, чтобы на крайней левой границе были только орлы, а на крайней правой — только решки. Конфигурация с наименьшей энергией должна интерполировать между двумя вакуумами. Эта граница раздела — доменная стенка.

\begin{figure}[htbp]
\centering
\includegraphics[width=0.9\textwidth]{lecture_06_figure_08.png}
\caption{Решётка орёл--решка, разделённая пунктирной доменной стенкой. Противоположные граничные условия вынуждают систему содержать границу раздела между двумя вырожденными упорядоченными состояниями.}
\lectureframe{01:30:35}
\end{figure}

\begin{classroomqa}
\audiencequestion{Чем спонтанное нарушение симметрии отличается от обычного явного нарушения симметрии?
\lecturetimestamp{01:28:36}}
\lecturerresponse{При явном нарушении локальный гамильтониан уже приписывает орлу и решке разные энергии, так что пары основных состояний, связанных симметрией, нет. При спонтанном нарушении уравнения сохраняют симметрию, вакуумы вырождены, и одно из них выбирается граничными условиями или сколь угодно малым возмущением. Возможность доменной стенки, интерполирующей между различными вакуумами, служит диагностическим признаком второй ситуации.}
\end{classroomqa}

Эта модель намеренно примитивна, но в ней содержится существенная логика, которая будет перенесена на поле Хиггса: симметричная теория, неустойчивая или вырожденная структура вакуума, выбранное основное состояние и физические следствия, определяемые этим выбором.

\section{К чему нас подводит лекция}

Таблица фермионов превращает группу Стандартной модели в набор физических направлений. Цветовое \(SU(3)\) действует по вертикали, переводя между красным, зелёным и синим; заряженные слабые преобразования действуют по горизонтали внутри кварковых и лептонных дублетов; фазовая симметрия даёт предварительный фактор \(U(1)\). Испускание калибровочных бозонов реализует эти групповые преобразования как взаимодействия.

Тогда слабость бета-распада объясняется не исключительно малым \(g_w\), а тяжёлым переносчиком. В координатном пространстве масса отсекает пропагатор на расстояниях, превышающих комптоновский масштаб \(\hbar/(m_Wc)\). В импульсном пространстве при ядерных импульсах это даёт
\begin{equation}
  \mathcal A_{\rm weak}
  \sim\frac{g_w^2}{k^2+m_W^2}
  \simeq\frac{g_w^2}{m_W^2}
\end{equation}
В старых низкоэнергетических экспериментах именно эта комбинация измерялась как константа Ферми.

Развёрнутое обсуждение в аудитории проясняет, что означает и чего не означает эта формула. Внутренние линии \(W\) виртуальны и находятся вне массовой оболочки; энергия--импульс всё же сохраняется. Достаточно жёсткий зонд может породить реальный \(W\), но тогда он создаёт другой процесс с иным конечным состоянием. Продукты столкновения не являются буквальным списком заранее существовавших составляющих. Температура изменяет число энергичных налетающих частиц, фотоны могут взаимодействовать через заряженные виртуальные частицы, а статистика Бозе отлична от динамической силы.

Наконец, \(SU(2)\) требует недостающего нейтрального калибровочного поля, связанного с \(Z\), а решётка орёл--решка даёт первую модель спонтанного нарушения симметрии. Следующий шаг — соединить эти две незавершённые линии: электрослабое смешивание и механизм Хиггса.
